\documentclass{article}
\usepackage[utf8]{inputenc}
\usepackage{amsmath}
\usepackage{gensymb}
\usepackage{setspace}
\usepackage{ragged2e}
\usepackage{enumerate}
\usepackage{amsthm}
\usepackage{graphicx}
\usepackage{subcaption}
\newtheorem{problem}{Problem}
\newtheorem{proposition}{Proposition}
\newtheorem*{solution}{Solution}
\usepackage{amssymb}
\usepackage{mathtools}
\usepackage{xcolor}
\usepackage{array}
\usepackage[T1]{fontenc}
\usepackage[margin=1in]{geometry}
\usepackage{lmodern}
\usepackage{tabularx}
\usepackage{pesmath}
\usepackage{fancyhdr}
\usepackage{graphicx}
\graphicspath{{images/}}

\newtheorem{lemma}{Lemma}


\pagestyle{fancy}
\fancyhf{}
\rhead{Group 1}
\chead{Final project}
\lhead{EN.553.740}
\rfoot{Page \thepage}





\usepackage{booktabs}
\usepackage{placeins}
\usepackage[hidelinks]{hyperref}
\setlength{\headheight}{14pt}
\newcommand{\TBD}[1]{\par\noindent\textcolor{blue!65!black}{\textbf{TBD:} #1}\par}
\newcommand{\placeholderfigure}[1]{%
  \fbox{\parbox[c][1.0in][c]{0.94\linewidth}{\centering\textbf{Figure placeholder}\\#1}}}

% MAIN REPORT BUDGET: 8 pages TOTAL, including title, abstract, figures and tables.
% Suggested allocation (not forced page breaks):
% title/abstract/introduction 1.25; related work 0.5; data 0.75;
% methods/theory 1.5; experiment design 0.75; results 2.0;
% discussion/conclusion 1.25. Total = 8 pages.
% References and appendices are outside this budget. Keep all essential
% contributions, strongest evidence and major conclusions BEFORE main-report-end.
% Do not rely on appendix review or reduce font/margins to accommodate extra text.

\title{\textbf{Interpretable Gene Programs or Foundation Models?}\\
\large Comparing Sparse Dictionary Learning and scGPT\\for Single-Cell Classification}
\author{Jonathan Ma \and Sijia Zhu \and Jiatong Gao \and Yunwei Chai}
\date{EN.553.740 Machine Learning 1\\Instructor: Dr. Sichen Yang\\\textit{Draft report scaffold}}

\begin{document}
\maketitle
\thispagestyle{fancy}

\begin{abstract}
\TBD{Write a concise abstract stating the question, controlled comparison,
dataset, main quantitative result, interpretability finding, and conclusion.
Compare processed expression, sparse dictionary codes, and frozen scGPT
embeddings using the same logistic-regression model family. Insert final results
only after model selection is frozen and the final evaluation is complete.}
\end{abstract}

\section{Introduction}
\TBD{Motivate the tradeoff between predictive utility and directly inspectable
gene programs in single-cell representation learning. State the research question:
what does the pretrained scGPT representation add beyond a sparse representation
learned from the target training data?}
\TBD{State the final contributions in two or three concrete points: a controlled
classification comparison, an analysis of sparse gene programs, and a
reproducible evaluation. Revise these to match the work actually completed.}

\section{Related Work}
\TBD{Briefly connect Lasso and sparse coding to dictionary learning
(Tibshirani; Olshausen and Field; Mairal et al.), gene-program analysis
(Kotliar et al.), scGPT (Cui et al.), and pretraining-scale evaluation
(DenAdel et al.). Add verified citations from the literature in jm.tex.
Explain how our signed dictionary differs from nonnegative factorization.}

\section{Data and Shared Evaluation Cohort}
\subsection{Dataset and provenance}
The project uses the Baron human pancreas dataset (GSE84133), obtained through
the \texttt{scRNAseq} function \texttt{BaronPancreasData}. The prepared dataset contains 8,569 cells,
20,125 genes, four donors, and 14 annotated cell types.
\TBD{Cite the original dataset and software; summarize count validation,
metadata preservation, and any relevant limitations of the annotations.}

\subsection{Exclusions and fixed splits}
Classes with fewer than 30 cells are excluded from the primary comparison,
removing 63 cells from the T-cell, Schwann, epsilon, and mast classes. The
remaining 8,506 cells span ten classes. A cell-type-stratified random split
with seed 42 assigns 5,954 cells to training, 1,277 to validation, and 1,275 to
test, approximating a 70/15/15 allocation. Donors can occur in multiple splits;
this evaluates new cells from the same donor pool, not unseen-donor transfer.
No additional cells were excluded by the scGPT panel check.
\TBD{Confirm these counts against the final run manifest. Explain the practical
rare-class cutoff and small remaining evaluation classes. Keep the split's
scientific limitation in the main report, not only in an appendix.}

\section{Methods and Theoretical Basis}
\subsection{Classical expression preprocessing}
The current baseline normalizes each cell to 10,000 total counts across the
original gene set, applies $\log(1+x)$, selects 2,000 genes by training-set
log-expression variance, and uses training-fitted gene-wise centering and
scaling. Identical fitted transforms are applied to validation and test cells.
\TBD{Confirm or revise this provisional recipe after coordination with the
classifier teammate. Specify that variance ranking is not dispersion-adjusted
HVG selection. Explain that scGPT has separate input processing.}

\subsection{Sparse dictionary learning}
Let $X\in\mathbb{R}^{n\times p}$ contain processed expression, with cells in
rows. We represent it by $ZD$, where $Z\in\mathbb{R}^{n\times K}$ contains
cell activities and $D\in\mathbb{R}^{K\times p}$ contains gene-program atoms:
\begin{equation}
\min_{Z,D}\;\frac12\|X-ZD\|_F^2+\lambda\|Z\|_1,
\qquad \|d_k\|_2\leq1\quad(k=1,\ldots,K).
\label{eq:dictionary}
\end{equation}
The atom constraint prevents reducing the sparsity penalty by rescaling the
factors. With $D$ fixed, each cell's code solves a convex Lasso problem; learning
both factors is nonconvex. Validation and test codes must use the learned
training dictionary without updating its atoms.
\TBD{Describe the final mini-batch solver, signed coefficients, stopping rule,
initialization, and selected $K$ and $\lambda$. Briefly explain the
reconstruction--sparsity tradeoff and limits of convergence claims. Keep this
core argument here; put detailed updates and derivations in the appendix.}

\subsection{Frozen scGPT embeddings}
The supplied comparison checkpoint is
\texttt{sctab\_seed0\_n1600000\_scgpt}: a small custom scGPT model pretrained on
1.6 million cells with sampling seed 0. It is not the original published large
scGPT checkpoint. Its frozen representation uses 128-dimensional,
L2-normalized CLS embeddings. The Baron input matches 1,879 of the fixed 2,000
panel genes; missing panel genes are zero-filled by the extractor.
\TBD{Cite scGPT, summarize the checkpoint choice made without Baron
classification scores, and describe binning and extraction settings. Verify
against the supplied metadata. Pretraining overlap remains unaudited unless a
subsequent audit resolves it; do not claim independent pretraining/test data.}

\subsection{Common logistic classifier}
\TBD{Describe three separately fitted L2-regularized multinomial logistic
classifiers: processed expression directly, dictionary codes, and scGPT
embeddings. Specify the common tuning procedure, class ordering, any
training-fitted representation scaling, and optimization settings. Do not
concatenate representations unless explicitly introducing a separate experiment.}

\section{Experimental Design}
\TBD{Describe validation-only selection and the agreed search budget. Record
revisions to the provisional dictionary grid after the classifier teammate's
initial validation results. The starter grid is $K\in\{16,32\}$ and
$\lambda\in\{1,5\}$, seed 42, with a ten-epoch ceiling; this is not yet the
final model-selection result. State how classifier regularization is tuned
consistently across representations.}
\TBD{Specify Macro-F1 as primary, plus balanced accuracy and multiclass log
loss. Explain cell-ID alignment, identical evaluation cells, and the rule that
final test results are inspected only after choices are frozen. If reporting
uncertainty or repeated runs, define the procedure and what variability it
captures; do not assume cells from the same donor are independent.}

\section{Results}
\subsection{Classification performance}
\TBD{Lead with the strongest quantitative comparison and answer the research
question. Keep validation selection separate from final held-out performance.
Do not fill this table with reconstruction metrics or smoke-test outputs.}
\begin{table}[htbp]
\centering
\caption{Final held-out classification results (not yet evaluated).
Higher Macro-F1 and balanced accuracy are better; lower log loss is better.}
\label{tab:classification}
\begin{tabular}{lccc}
\toprule
Representation & Macro-F1 & Balanced accuracy & Log loss\\
\midrule
Processed expression & TBD & TBD & TBD\\
Sparse dictionary codes & TBD & TBD & TBD\\
Frozen scGPT embeddings & TBD & TBD & TBD\\
\bottomrule
\end{tabular}
\end{table}

\subsection{Predictive utility and sparsity}
\begin{figure}[htbp]
\centering
\placeholderfigure{Validation performance versus dictionary size/penalty,\\with code sparsity indicated.}
\caption{TBD: show the principal tradeoff supporting the selected dictionary.
Specify the evaluation split and uncertainty, if any.}
\label{fig:tradeoff}
\end{figure}
\TBD{Discuss whether compression and sparse activity preserve useful cell-type
information. Report reconstruction as a diagnostic rather than evidence of
classification success. Summarize convergence and stability only where tested.}

\subsection{Gene-program interpretation}
\begin{figure}[htbp]
\centering
\placeholderfigure{Selected atoms: high-weight genes and\\cell-type activation patterns.}
\caption{TBD: strongest interpretable gene-program evidence. Distinguish positive
and negative weights and describe the selection criterion for displayed atoms.}
\label{fig:programs}
\end{figure}
\TBD{Discuss a small number of informative programs with biological support
where available. Account for atom sign/permutation ambiguity. Avoid equating
associations with causal mechanisms or claiming every atom is a unique pathway.}

\section{Discussion and Limitations}
\TBD{Explain the performance--interpretability tradeoff supported by the final
results, including whether scGPT adds linearly accessible predictive information.
Discuss why the expression baseline matters. Avoid conclusions based solely on
model size or reconstruction.}
\TBD{Keep major limits in the main text: shared donors across splits, excluded
rare classes, small per-class evaluation samples, missing panel genes,
custom-checkpoint scope, unaudited pretraining overlap, and validation-search
limitations. Add computational or stability limitations observed in final runs.}

\section{Conclusion}
\TBD{Give the direct answer to the research question and the principal evidence
in a short paragraph. State what remains unresolved. Do not introduce a new
experiment or essential qualification only in the appendix.}

% This marker must remain BEFORE references and appendices.
\FloatBarrier
\label{main-report-end}
\typeout{MAIN REPORT ENDS ON PAGE \thepage; LIMIT IS 8 PAGES.}
\clearpage

\section*{References}
% Replace this placeholder with a bibliography, or a thebibliography environment.
% References do not count toward the eight-page main-report limit.
\TBD{Insert the verified, cited bibliography. Start from jm.tex, and add the
Baron dataset and implementation/software references actually used. Ensure every
in-text citation resolves and every claim is supported by the relevant source.}

\clearpage
\appendix
\section{Supplementary Theory}
\TBD{Detailed sparse-coding derivation, proximal/alternating updates, atom-norm
constraint reasoning, and additional optimization discussion. The main text
must retain the objective and key theoretical insights.}

\section{Implementation and Reproducibility}
\TBD{Commands, package versions, checkpoint/data/manifest hashes, artifact
schemas, seeds, batch sizes, precision, and hardware. Record the final code
revision and access instructions without including private credentials.}

\section{Supplementary Experiments and Figures}
\TBD{Complete candidate-grid results, per-class metrics, confusion matrices,
additional atoms, convergence diagnostics, and seed sensitivity where available.
Only supplementary evidence belongs here; the strongest results remain in the
main report. Clearly label analyses that were not performed.}

\section{Team Contributions}
\TBD{Summarize each member's actual contributions to preprocessing, scGPT,
dictionary learning, classification, verification, and writing.}

\end{document}
